\documentclass[11pt,a4paper]{article}
\usepackage[margin=22mm]{geometry}
\usepackage{amsmath,amssymb,amsthm,mathtools,booktabs,microtype}
\usepackage[hidelinks]{hyperref}
\newtheorem{proposition}{Proposition}
\newcommand{\norm}[1]{\lVert#1\rVert}
\title{GR-03: Shifted Volterra product integration\\\large Linear and quadratic rules for the two-kernel Gr\"onwall framework}
\author{Research working notes}
\date{26 November 2026 --- revision 01}
\begin{document}
\maketitle
\noindent\textbf{Revision 01 (AUD-02 closure).}
The independent analytical review verdict is \emph{ACCEPT WITH LIMITATIONS}.
The results below are conditional on their stated hypotheses. Finite-grid
diagnostics are evidence for implementations, not proofs of the analytical
theorems. The numerical tables retain the original diagnostic values.
Lean certificates cover the explicitly identified finite algebraic subclaims,
not complete fractional PDE arguments. Original manuscripts and execution
evidence are preserved separately.

\noindent\textbf{Scope.} This note continues GR-01 and GR-02, independently of the frozen AUD-01 checkpoint. It derives two causal uniform-grid memory rules and proves the absolute row-tail bounds needed by the complementary-kernel comparison. The constructions are polynomial product integration; no novelty claim is made. All orders are fixed in $(0,1)$; the results do not require $\beta\le\alpha$ except where explicitly imposed in examples. No new production algorithm, nonuniform-grid theorem or full convergence theorem is claimed.

\section*{Review qualifications}
These explicit coefficients and constants are for a uniform mesh and fixed
orders. GR-06 rederives the nonuniform stencil geometry and its
grading-dependent absolute-basis constant; the value $5/4$ must not be
transferred unchanged to an arbitrary backward nonuniform stencil.

The rule approximates normalized $I^r$; the equation includes
$\lambda\Gamma(r)$ exactly once. It is causal implicit: completed forward
stencils use levels at most $n$, and the final quadratic memory stencil is
backward. At startup both memory rules are linear. The final Caputo panel,
the final memory panel, and the shifted diffusion average have distinct roles.

The exact negative witness is the \emph{normalized factor}
$W^Q_{20}/[t_*^r/\Gamma(1+r)]=-17/231$, not the dimensional weight itself.
All later comparisons use $|W^Q|$, including its initial column.
Sharpness of the startup memory error, post-startup memory error and full
solution error are different statements. A memory estimate alone neither
proves nor refutes a PDE convergence theorem.

Nonpolynomial nodal coefficients in spatial applications require continuous
representatives with pointwise bounds at the nodes; an essential bound on an
arbitrarily modified equivalence class is insufficient.
The Lean certificates check finite moment collection, the uniform basis
bound and the rational witness. Gamma/beta integration, consistency and the
PDE energy application remain analytical.

\section{Normalized memory and common moments}
Write $r=1-\beta$, $\theta=\alpha/2$, $\sigma=1-\theta$, $t_k=k\tau$, $t_*=t_{n-1}+\sigma\tau$, and $T=N\tau$. The original unnormalized memory $\mathcal J$ is $\Gamma(r)I^r$, where
\[
 I^r v(t)=\frac1{\Gamma(r)}\int_0^t(t-s)^{r-1}v(s)\,ds.
\]
We discretize $I^r v(t_*)$; the PDE coefficient multiplying the rule is therefore $\lambda\Gamma(r)K(t_*)$. On panel $k$ use $x=(s-t_{k-1})/\tau$, $d=n-k+\sigma$, and $h=1$ for $k<n$, $h=\sigma$ for $k=n$. Define the moments
\begin{align*}
 M_j(d,h)&=\frac{\tau^r}{\Gamma(r)}\int_0^h x^j(d-x)^{r-1}\,dx\\
 &=\frac{\tau^rd^{r+j}}{\Gamma(r)}B_{h/d}(j+1,r),\\
 &=\frac{\tau^r}{\Gamma(r)}\sum_{i=0}^j(-1)^i\binom ji d^{j-i}
       \frac{d^{r+i}-(d-h)^{r+i}}{r+i},
\end{align*}
Here $j=0,1,2$ and $B_z$ is the incomplete beta function. The last expression follows by expanding $(d-y)^j$ with $y=d-x$. It is algebraically exact but may suffer cancellation at large lag. The diagnostic uses the regularized incomplete beta function and a separate adaptive quadrature oracle; neither establishes a floating-point stability theorem.
\section{Linear and quadratic assembly}
The linear rule PI-L integrates the interpolant $(1-x)U^{k-1}+xU^k$ on every panel, including the last partial panel. Add $M_0-M_1$ and $M_1$ to columns $k-1,k$, respectively. On the final panel $M_0=(\sigma\tau)^r/\Gamma(1+r)$ and $M_1=\sigma M_0/(1+r)$. Thus its weights are not the shifted evaluation weights $\theta,\sigma$.

PI-Q uses the forward stencil $k-1,k,k+1$ on all completed panels and the backward stencil $n-2,n-1,n$ on the final partial panel ($n\ge2$). At $n=1$, PI-Q uses PI-L. In both cases all data are at levels at most $n$. The following table lists the local basis polynomials and exact contributions to the assembled row.
\begin{center}\small
\begin{tabular}{lll}
\toprule
Panel/stencil & Basis & Contribution\\\midrule
Forward: $k-1$ & $(x-1)(x-2)/2$ & $M_0-3M_1/2+M_2/2$\\
$k$ & $x(2-x)$ & $2M_1-M_2$\\
$k+1$ & $x(x-1)/2$ & $(M_2-M_1)/2$\\\midrule
Backward: $n-2$ & $x(x-1)/2$ & $(M_2-M_1)/2$\\
$n-1$ & $1-x^2$ & $M_0-M_2$\\
$n$ & $x(x+1)/2$ & $(M_2+M_1)/2$\\\bottomrule
\end{tabular}
\end{center}
Let $Q_X^n[U]=\sum_{k=0}^n W^X_{nk}U^k$ for $X=L,Q$. Both rules integrate constants and linear polynomials exactly; PI-Q integrates quadratic polynomials exactly for $n\ge2$.
\begin{proposition}[Mass and sign]
For both rules $\sum_{k=0}^nW^X_{nk}=t_*^r/\Gamma(1+r)$. The PI-L entries are nonnegative. PI-Q can have negative assembled entries, but
\[
 \sum_{k=0}^n|W^X_{nk}|\le L_X\frac{t_*^r}{\Gamma(1+r)},
 \qquad L_L=1,\quad L_Q=5/4.
\]
\end{proposition}
\begin{proof}
Partition of unity gives the signed mass. The linear basis is nonnegative. On either quadratic stencil inside its integration interval the sum of absolute basis functions is $1+x(1-x)\le5/4$. Integrate against the positive kernel and apply the triangle inequality to assembly. The startup row obeys the stronger linear bound.
\end{proof}
An exact signed-weight witness is obtained at $n=2$: both quadratic pieces are the same interpolating polynomial on the nodes $0,\tau,2\tau$. Put $D=t_*/\tau=1+\sigma$. The initial-data basis is $(z-1)(z-2)/2$ with $z=s/\tau$. Integrating its three monomials gives
\[
 W^Q_{20}=\frac{t_*^r}{\Gamma(1+r)}\left[1-\frac{3D}{2(1+r)}+\frac{D^2}{(1+r)(2+r)}\right].
\]
At $\alpha=\beta=9/10$ the bracket is exactly $-17/231$. Hence positivity of the assembled quadratic rule cannot be used in a general energy proof.
\section{Smooth consistency and sharp startup order}
All statements below hold for Hilbert-space valued data; scalar interpolation bounds extend by duality. The standard interpolation remainder yields
\begin{align*}
 \norm{\Pi_1u-u}_{\infty}&\le \tau^2\norm{u''}_{\infty}/8,\\
 \norm{\Pi_2u-u}_{\infty}&\le \tau^3\norm{u'''}_{\infty}/12.
\end{align*}
For the second bound the nodal product has magnitude at most $\tau^3/2$ on either stencil, since $x(1-x)\le1/4$ and the remaining factor is at most 2. Multiplying by the exact positive kernel mass proves
\begin{proposition}[Consistency]
For $u\in C^2([0,T];H)$,
\[
 \norm{Q_L^n[u]-I^ru(t_*)}\le\frac{\tau^2 t_*^r}{8\Gamma(1+r)}\norm{u''}_{\infty}.
\]
For $u\in C^3([0,T];H)$, PI-Q satisfies
\[
 \norm{Q_Q^n[u]-I^ru(t_*)}\le
 \begin{cases}
 \sigma^r\tau^{2+r}\norm{u''}_{\infty}/(8\Gamma(1+r)),&n=1,\\
 \tau^3t_*^r\norm{u'''}_{\infty}/(12\Gamma(1+r)),&n\ge2.
 \end{cases}
\]
In particular the uniform smooth consistency order of the linearly started PI-Q is $2+r=3-\beta$, not 3.
\end{proposition}
The startup bound is sharp: for $u(t)=t^2$, the first-step error equals
\[
 \frac{\tau^{2+r}}{\Gamma(r)}\int_0^\sigma x(1-x)(\sigma-x)^{r-1}\,dx>0,
\]
and all subsequent quadratic-rule errors vanish. The PI-L order is sharp at $n=N$ for $u(t)=t^2$: the error is nonnegative on every panel, and on each completed panel its kernel-weighted integral is at least $t_*^{r-1}\tau^3/(6\Gamma(r))$. Summing $N-1$ such contributions gives a positive constant times $\tau^2$ for fixed $T$ and sufficiently fine grids.

The post-startup PI-Q order is also sharp for $u=t^3$ at $n=N$. On completed forward panels $u-\Pi_2u=\tau^3 x(x-1)(x-2)\ge0$. Their total integrated error has a lower bound $(N-1)t_*^{r-1}\tau^4/(4\Gamma(r))$, using $\int_0^1 x(x-1)(x-2)\,dx=1/4$. The opposite-sign final backward-panel contribution is $O(\tau^{3+r})$, so it cannot cancel the $\tau^3$ leading lower bound for sufficiently small $\tau$. These are memory consistency results, not solution convergence statements. The shifted diffusion approximation still has a second-order interpolation error.
\section{Localized row-tail bounds}
For $X=L,Q$, define the nonnegative majorant $H^X_{nk}=|W^X_{nk}|$, retaining the initial-data column $k=0$. Let $1\le m\le n$ and $\ell=t_n-t_{m-1}$. In the linear rule a basis with index $k\ge m$ has support only for $s\ge t_{m-1}$. In the quadratic rule the forward stencil allows one additional earlier panel, but none before $t_{\max(0,m-2)}$. The final backward stencil respects the latter restriction as well. Since $\ell\ge\tau$, the length from this earliest point to $t_*$ is at most $\ell+\tau\le2\ell$. The same pointwise absolute-basis bounds as above then yield
\begin{proposition}[Localized memory mass]
\[
 \sum_{k=m}^nH^X_{nk}\le
 \frac{C_X}{\Gamma(1+r)}(t_n-t_{m-1})^r,
 \qquad C_L=1,\quad C_Q=\tfrac54\,2^r.
\]
\end{proposition}
The factor $2^r$ is a simple uniform-grid estimate of the one-panel stencil extension; it is not asserted optimal. In particular, a linear-rule tail bound must not be applied unchanged to the quadratic stencil.

Let $P$ be the nonnegative GR-02 complementary matrix, indexed by $1,\ldots,N$, satisfying
\[
 \sum_{j=m}^n P_{nj}\le d_\alpha^{-1}(t_n-t_{m-1})^\alpha,
 \qquad d_\alpha=(2\Gamma(1-\alpha))^{-1}.
\]
In products with $P$ restrict $H^X$ to columns $1,\ldots,N$ unless the initial column is explicitly included.
\begin{proposition}[Two fractional powers]
\[
 \sum_{k=m}^n(PH^X)_{nk}\le
 \frac{C_X}{d_\alpha\Gamma(1+r)}(t_n-t_{m-1})^{\alpha+r}.
\]
\end{proposition}
\begin{proof}
Reverse the finite sums: the left side equals $\sum_{j=m}^nP_{nj}\sum_{k=m}^jH^X_{jk}$. Use the previous tail bound, $t_j-t_{m-1}\le\ell$, and the tail bound for $P$. No discrete semigroup or commutativity identity is used.
\end{proof}
\section{Bridge to GR-04: shifted energy and initial data}
Consider the shifted equation with $A$ a nonnegative self-adjoint diffusion operator, $K_*$ a bounded operator, $v_*^n=\sigma U^n+\theta U^{n-1}$, and memory $\lambda\Gamma(r)K_*Q_X^n[U]$. For an existing solution in the operator domain, put $E^k=\norm{U^k}^2$, $\mu=|\lambda|\Gamma(r)\sup_n\norm{K_*}$ and $F_n=\norm{f_*^n}^2$. Testing with $2v_*^n$, using GR-02's energy inequality, convexity, and Young's inequality, gives for $\varepsilon>0$
\begin{equation*}
 (\mathcal D_\tau^\alpha E)^{n-\theta}\le
 b(\sigma E^n+\theta E^{n-1})+\mu\sum_{k=0}^nH^X_{nk}E^k+\varepsilon^{-1}F_n,
 \quad b=\varepsilon+\mu L_X T^r/\Gamma(1+r).
\end{equation*}
Indeed $2\sqrt{E^k}\norm{v_*}\le E^k+\norm{v_*}^2$ and $\norm{v_*}^2\le\sigma E^n+\theta E^{n-1}$. This step needs only absolute memory mass, not positive $W^Q$.

Define the rectangular shift matrix $S_{nk}=\sigma\mathbf1_{k=n}+\theta\mathbf1_{k=n-1}$ for $0\le k\le N$; denote the initial column by $S_0$ and the remaining square matrix by $S_+$. Use analogous $H_0,H_+$ notation. Multiplication by $P$ gives, for $E_+=(E^1,\ldots,E^N)^\top$,
\begin{align*}
 E_+&\le a+B E_+,\qquad B=bPS_++\mu PH_+,\\
 a&=E^0(\mathbf1+bPS_0+\mu PH_0)+\varepsilon^{-1}PF.
\end{align*}
Since $\sum_{k=m}^j(S_+)_{jk}\le1$ and is zero for $j<m$, the row tail of $PS_+$ is at most $d_\alpha^{-1}\ell^\alpha$. The previous proposition bounds $PH_+$. Thus the nonnegative matrix $B$ has tail at most
\[
 d_\alpha^{-1}\left(b\ell^\alpha+\frac{\mu C_X}{\Gamma(1+r)}\ell^{\alpha+r}\right).
\]
The initial column contributes a bounded data term:\[PS_0\le d_\alpha^{-1}T^\alpha\mathbf1,\qquad PH_0\le \frac{L_X T^{\alpha+r}}{d_\alpha\Gamma(1+r)}\mathbf1.\] The GR-01 block argument therefore yields a conditional finite-time energy bound for any existing solution of either scheme for sufficiently small $\tau$. This is the exact algebraic bridge, not yet a full GR-04 analysis: per-step solvability, a more quantitative discrete resolvent, error propagation in sharper norms, and regularity assumptions for full convergence are reserved for subsequent work.
\section{Numerical diagnostic and limitations}
The standalone \texttt{verify\_gr03.py} uses incomplete-beta moments, independent adaptive quadrature of Lagrange basis functions (with an algebraic endpoint weight on the final panel), and a separately assembled GR-02 complementary matrix. Six $(\alpha,\beta)$ pairs, three grid sizes and two rules give 36 configurations. All revised assertions passed. Maximum discrepancies: direct/assembled weights $6.661\times10^{-16}$, constant mass $3.331\times10^{-16}$, exact polynomials $5.551\times10^{-16}$. Maximum sampled normalized tail ratios are $0.99936760$ for $H$, $0.35917033$ for $PH$. The tests detect negative assembled quadratic entries.
\begin{center}\small
\begin{tabular}{ccrrr}
\toprule
$\alpha$ & $\beta$ & PI-L $t^2$ & PI-Q $t^2$ startup & PI-Q $t^3$, $n\ge2$\\\midrule
0.50&0.25&2.000970&2.750000&2.970834\\
0.75&0.50&2.005973&2.500000&2.921839\\
0.95&0.90&2.019041&2.100000&3.100000\\\bottomrule
\end{tabular}
\end{center}
Rates are last-doubling observations ($N=128$ to 256), not full PDE rates. The last quadratic rate is pre-asymptotic, with possible cancellation of signed panel contributions. The initial diagnostic failed during the convergence loop because a global basis-array name was overwritten by an assembled matrix; the original source and partial output are retained. Only variable names were corrected, and no assertion was weakened. Python packages: NumPy 2.3.5, SciPy 1.17.0, mpmath 1.3.0. The later independent review and finite Lean coverage do not turn these historical diagnostics into analytical proofs.

\textbf{GR-04 target:} establish existence and mesh-uniform stability of the combined scheme with a fully quantified discrete two-kernel comparison; then use that result for convergence and initial-layer analysis. PI-L remains the natural primary scheme, with PI-Q a controlled signed-weight alternative.
\end{document}
